\documentclass[10pt,a4paper]{amsart}
\usepackage{amsmath,amssymb,amsthm,graphicx}
% Retained preamble for the Stacks Project chapter snapshot in this directory.
% The chapter's own preamble.tex is a copy of the shared Stacks Project preamble; the block
% below is sliced MECHANICALLY, line for line, from it. Nothing was added, removed or
% redefined; the wrapper class, the xr-hyper/hyperref wiring and the externaldocument list
% are omitted because the assembler supplies the document class.
%
% lines taken from preamble.tex: 3, 5, 9, 11, 12, 156, 157, 158, ... (51 lines total)
%
% package coverage restored from the chapter's own preamble (2026-10-13): xy, xy-2cell

\usepackage[all]{xy}
\xyoption{2cell}
\UseAllTwocells

\newenvironment{reference}{\comment}{\endcomment}
\newenvironment{slogan}{\comment}{\endcomment}
\newenvironment{history}{\comment}{\endcomment}
\theoremstyle{plain}
\newtheorem{theorem}[subsection]{Theorem}
\newtheorem{proposition}[subsection]{Proposition}
\newtheorem{lemma}[subsection]{Lemma}
\theoremstyle{definition}
\newtheorem{definition}[subsection]{Definition}
\newtheorem{example}[subsection]{Example}
\newtheorem{exercise}[subsection]{Exercise}
\newtheorem{situation}[subsection]{Situation}
\theoremstyle{remark}
\newtheorem{remark}[subsection]{Remark}
\newtheorem{remarks}[subsection]{Remarks}
\numberwithin{equation}{subsection}
\def\lim{\mathop{\mathrm{lim}}\nolimits}
\def\colim{\mathop{\mathrm{colim}}\nolimits}
\def\Spec{\mathop{\mathrm{Spec}}}
\def\Hom{\mathop{\mathrm{Hom}}\nolimits}
\def\Ext{\mathop{\mathrm{Ext}}\nolimits}
\def\SheafHom{\mathop{\mathcal{H}\!\mathit{om}}\nolimits}
\def\SheafExt{\mathop{\mathcal{E}\!\mathit{xt}}\nolimits}
\def\Sch{\mathit{Sch}}
\def\Mor{\mathop{\mathrm{Mor}}\nolimits}
\def\Ob{\mathop{\mathrm{Ob}}\nolimits}
\def\Sh{\mathop{\mathit{Sh}}\nolimits}
\def\NL{\mathop{N\!L}\nolimits}
\def\CH{\mathop{\mathrm{CH}}\nolimits}
\def\proetale{{pro\text{-}\acute{e}tale}}
\def\etale{{\acute{e}tale}}
\def\QCoh{\mathit{QCoh}}
\def\Ker{\mathop{\mathrm{Ker}}}
\def\Im{\mathop{\mathrm{Im}}}
\def\Coker{\mathop{\mathrm{Coker}}}
\def\Coim{\mathop{\mathrm{Coim}}}
\DeclareMathSymbol{\boxtimes}{\mathbin}{AMSa}{"02}
\def\QCohstack{\mathcal{QC}\!\mathit{oh}}
\def\Cohstack{\mathcal{C}\!\mathit{oh}}
\def\Spacesstack{\mathcal{S}\!\mathit{paces}}
\def\Quotfunctor{\mathrm{Quot}}
\def\Hilbfunctor{\mathrm{Hilb}}
\def\Curvesstack{\mathcal{C}\!\mathit{urves}}
\def\Polarizedstack{\mathcal{P}\!\mathit{olarized}}
\def\Complexesstack{\mathcal{C}\!\mathit{omplexes}}
\def\Pic{\mathop{\mathrm{Pic}}\nolimits}
\def\Picardstack{\mathcal{P}\!\mathit{ic}}
\def\Picardfunctor{\mathrm{Pic}}
\def\Deformationcategory{\mathcal{D}\!\mathit{ef}}
\usepackage{cleveref}
\crefname{originallemma}{original lemma}{original lemmas}
\Crefname{originallemma}{Original lemma}{Original lemmas}
\makeatletter
\newlabel{lemma-dgm-abelian}{{0FRJ}{source}{}{sourcepointer-1}{}}
\newlabel{lemma-dgm-abelian@cref}{{[originallemma][0][]0FRJ}{[1][0][]source}{}{}{}}
\makeatother
\crefname{originallemma}{original lemma}{original lemmas}
\Crefname{originallemma}{Original lemma}{Original lemmas}
\makeatletter
\newlabel{lemma-hom-derived}{{0FT8}{source}{}{sourcepointer-2}{}}
\newlabel{lemma-hom-derived@cref}{{[originallemma][0][]0FT8}{[1][0][]source}{}{}{}}
\makeatother
\crefname{originallemma}{original lemma}{original lemmas}
\Crefname{originallemma}{Original lemma}{Original lemmas}
\makeatletter
\newlabel{lemma-homotopy-direct-sums}{{0FS2}{source}{}{sourcepointer-3}{}}
\newlabel{lemma-homotopy-direct-sums@cref}{{[originallemma][0][]0FS2}{[1][0][]source}{}{}{}}
\makeatother
\crefname{originallemma}{original lemma}{original lemmas}
\Crefname{originallemma}{Original lemma}{Original lemmas}
\makeatletter
\newlabel{lemma-product-K-injective}{{0FSV}{source}{}{sourcepointer-4}{}}
\newlabel{lemma-product-K-injective@cref}{{[originallemma][0][]0FSV}{[1][0][]source}{}{}{}}
\makeatother
\crefname{originallemma}{original lemma}{original lemmas}
\Crefname{originallemma}{Original lemma}{Original lemmas}
\makeatletter
\newlabel{lemma-product-graded-injective}{{0FSR}{source}{}{sourcepointer-5}{}}
\newlabel{lemma-product-graded-injective@cref}{{[originallemma][0][]0FSR}{[1][0][]source}{}{}{}}
\makeatother
\crefname{originallemma}{original lemma}{original lemmas}
\Crefname{originallemma}{Original lemma}{Original lemmas}
\makeatletter
\newlabel{theorem-qis-into-dg-injective}{{0FT0}{source}{}{sourcepointer-6}{}}
\newlabel{theorem-qis-into-dg-injective@cref}{{[originallemma][0][]0FT0}{[1][0][]source}{}{}{}}
\makeatother
\title{Stacks, Differential Graded Sheaves: lemma-derived-products}
\author{The Stacks Project Authors}
\date{stacks-project:sdga.tex:local-snapshot-20261012; GNU FDL 1.2 or later}
\begin{document}
\maketitle
\noindent\textit{Editorial note.} Statement span is lines 4551-4563 of sources/stacks-sdga/pinned-original-source.tex; proof span is lines 4565-4619, the complete author proof (53 body lines).\par\medskip
\noindent\textit{Editorial note.} Navigation stubs: 6 cited-but-not-included labels are supplied as non-load-bearing pointers with their pinned tags where the snapshot tag index has them.\par\medskip
\noindent\textit{Editorial note.} extra\_packages is empty on purpose: the retained preamble already loads xy with options and a second declaration clashes in the assembler wrapper.\par\medskip
\noindent\textit{Editorial note.} Text is retained verbatim, including the author's own typos and placeholders.\par\medskip
\noindent\textit{Editorial note.} Queue index 10 of queue\_sdga.json (pb=2234, stmt 4551-4563, proof 4565-4619).\par\medskip
\noindent\textit{Editorial note.} Cross-chapter same-label warning (NOT an occupancy conflict, Chief ruling 2026-10-13 evening): label 'lemma-derived-products' also exists in stacks-dga — a different proposition in a different chapter file.\par\medskip
\noindent\textit{Allowed standard prerequisites (editorial).}\begin{itemize}
\item Explicit prerequisite boundary. The cited inputs are prerequisites; the argument in the retained span is not replaced by any cited result.
\end{itemize}\medskip
\ifdefined\hypertarget\hypertarget{sourcepointer-1}{}\fi
\noindent\textit{Source reference.} Original lemma 0FRJ: sources/stacks-sdga/pinned-original-source.tex, source lines 1444--1455. Cited by the retained statement or proof but not part of the collected spans; supplied only as a navigation pointer. The argument does not depend on its content.\par\medskip
\ifdefined\hypertarget\hypertarget{sourcepointer-2}{}\fi
\noindent\textit{Source reference.} Original lemma 0FT8: sources/stacks-sdga/pinned-original-source.tex, source lines 4508--4520. Cited by the retained statement or proof but not part of the collected spans; supplied only as a navigation pointer. The argument does not depend on its content.\par\medskip
\ifdefined\hypertarget\hypertarget{sourcepointer-3}{}\fi
\noindent\textit{Source reference.} Original lemma 0FS2: sources/stacks-sdga/pinned-original-source.tex, source lines 2535--2541. Cited by the retained statement or proof but not part of the collected spans; supplied only as a navigation pointer. The argument does not depend on its content.\par\medskip
\ifdefined\hypertarget\hypertarget{sourcepointer-4}{}\fi
\noindent\textit{Source reference.} Original lemma 0FSV: sources/stacks-sdga/pinned-original-source.tex, source lines 3968--3976. Cited by the retained statement or proof but not part of the collected spans; supplied only as a navigation pointer. The argument does not depend on its content.\par\medskip
\ifdefined\hypertarget\hypertarget{sourcepointer-5}{}\fi
\noindent\textit{Source reference.} Original lemma 0FSR: sources/stacks-sdga/pinned-original-source.tex, source lines 3776--3784. Cited by the retained statement or proof but not part of the collected spans; supplied only as a navigation pointer. The argument does not depend on its content.\par\medskip
\ifdefined\hypertarget\hypertarget{sourcepointer-6}{}\fi
\noindent\textit{Source reference.} Original lemma 0FT0: sources/stacks-sdga/pinned-original-source.tex, source lines 4233--4242. Cited by the retained statement or proof but not part of the collected spans; supplied only as a navigation pointer. The argument does not depend on its content.\par\medskip
\begin{lemma}
\label{lemma-derived-products}
Let $(\mathcal{C}, \mathcal{O})$ be a ringed site. Let
$(\mathcal{A}, \text{d})$ be a sheaf of differential graded algebras
on $(\mathcal{C}, \mathcal{O})$. Then
\begin{enumerate}
\item $D(\mathcal{A}, \text{d})$ has both direct sums and products,
\item direct sums are obtained by taking direct sums of differential graded
$\mathcal{A}$-modules,
\item products are obtained by taking products of
K-injective differential graded modules.
\end{enumerate}
\end{lemma}

\begin{proof}
We will use that $\textit{Mod}(\mathcal{A}, \text{d})$ is an abelian category
with arbitrary direct sums and products, and that these give rise
to direct sums and products in $K(\textit{Mod}(\mathcal{A}, \text{d}))$.
See Lemmas \ref{lemma-dgm-abelian} and \ref{lemma-homotopy-direct-sums}.

\medskip\noindent
Let $\mathcal{M}_j$ be a family of differential graded $\mathcal{A}$-modules.
Consider the direct sum $\mathcal{M} = \bigoplus \mathcal{M}_j$
as a differential graded $\mathcal{A}$-module.
For a differential graded $\mathcal{A}$-module
$\mathcal{N}$ choose a quasi-isomorphism
$\mathcal{N} \to \mathcal{I}$ where
$\mathcal{I}$ is graded injective and K-injective as a
differential graded $\mathcal{A}$-module. See
Theorem \ref{theorem-qis-into-dg-injective}.
Using Lemma \ref{lemma-hom-derived} we have
\begin{align*}
\Hom_{D(\mathcal{A}, \text{d})}(\mathcal{M}, \mathcal{N})
& =
\Hom_{K(\mathcal{A}, \text{d})}(\mathcal{M}, \mathcal{I}) \\
& =
\prod \Hom_{K(\mathcal{A}, \text{d})}(\mathcal{M}_j, \mathcal{I}) \\
& =
\prod \Hom_{D(\mathcal{A}, \text{d})}(\mathcal{M}_j, \mathcal{I})
\end{align*}
whence the existence of direct sums in $D(A, \text{d})$ as given in
part (2) of the lemma.

\medskip\noindent
Let $\mathcal{M}_j$ be a family of differential graded $\mathcal{A}$-modules.
For each $j$ choose a quasi-isomorphism
$\mathcal{M} \to \mathcal{I}_j$ where
$\mathcal{I}_j$ is graded injective and K-injective as a
differential graded $\mathcal{A}$-module.
Consider the product $\mathcal{I} = \prod \mathcal{I}_j$
of differential graded $\mathcal{A}$-modules.
By Lemmas \ref{lemma-product-K-injective} and
\ref{lemma-product-graded-injective} we see that
$\mathcal{I}$ is graded injective and K-injective as a
differential graded $\mathcal{A}$-module.
For a differential graded $\mathcal{A}$-module
$\mathcal{N}$ using Lemma \ref{lemma-hom-derived} we have
\begin{align*}
\Hom_{D(\mathcal{A}, \text{d})}(\mathcal{N}, \mathcal{I})
& =
\Hom_{K(\mathcal{A}, \text{d})}(\mathcal{N}, \mathcal{I}) \\
& =
\prod \Hom_{K(\mathcal{A}, \text{d})}(\mathcal{N}, \mathcal{I}_j) \\
& =
\prod \Hom_{D(\mathcal{A}, \text{d})}(\mathcal{N}, \mathcal{M}_j)
\end{align*}
whence the existence of products in $D(\mathcal{A}, \text{d})$ as given in
part (3) of the lemma.
\end{proof}

\end{document}
